% --- Kata Pengantar ---
%
%     Perintah \KataPengantar digunakan agar dapat 
%     diterjemah saat mengganti bahasa ke Bahasa Indonesia 
%     atau Bahasa Inggris
%
%     \KataPengantar
%     Kata Pengantar <--> Preface
%
\chapter*{\KataPengantar}
\addcontentsline{toc}{chapter}{\KataPengantar}

%=========================================================%
%               TULIS KATA PENGANTAR DI SINI              %
% Jangan Lupa untuk Menghapus Contoh Tulisan di Bawah Ini %
%                                                         %
% [TIP] Sebaiknya jangan dibiasakan menulis gelar atau    %
% jabatan dari nama orang yang dicantumkan di Kata        %
% Pengantar. Penulisan yang tanpa mengurangi rasa hormat  %
% adalah dengan menyebutkan perannya.                     %
%                                                         %
% Contohnya "Gustin Galunggung, selaku profesor yang      %
% mengampu mata kuliah ... dan memberi bimbingan ..."     %
%=========================================================%

Puji syukur kehadirat Tuhan Yang Maha Esa atas rahmat dan karunia-Nya, sehingga penulis dapat menyelesaikan makalah dengan judul ``\judul'' ini dengan baik dan tepat waktu. Makalah ini disusun untuk memenuhi tugas mata kuliah.

Penyusunan makalah ini tidak lepas dari bimbingan dan bantuan dari berbagai pihak. Oleh karena itu, pada kesempatan ini penulis ingin mengucapkan terima kasih yang sebesar-besarnya kepada:

\begin{enumerate}[]
    \item \namaPengampuPolos, selaku tutor pengampu mata kuliah yang telah memberikan materi dan arahan selama proses penyusunan makalah ini.
    \item Orang tua dan keluarga yang selalu memberikan dukungan moral dan material.
    \item Teman-teman komunitas Prodi \prodi\ yang telah bekerja sama,  memberi saran, dan berbagi ilmu.
    \item Semua pihak yang tidak dapat penulis sebutkan satu per satu, atas kontribusi dan dukungannya.
\end{enumerate}

Penulis menyadari bahwa makalah ini masih jauh dari sempurna. Oleh karena itu, kritik dan saran yang membangun dari pembaca sangat penulis harapkan demi perbaikan di masa mendatang. Semoga makalah ini dapat memberikan manfaat serta wawasan bagi pembaca.

% --- Tempat, Tanggal, dan Penulis ---
\vfill
\hfill\begin{tabular}{@{}c@{}} 
    \daerahMahasiswa, \tanggalLengkap \vspace{.2cm} \\
    \vspace*{2.1cm} \\ % <-- Setinggi Meterai Rp10.000
    \namaMahasiswa
\end{tabular}
\vfill